\chapter{Anexos}


\section{Criterios SMART de los objetivos específicos}
\label{sec:anexo-smart}

Se desarrollan aquí los cinco criterios SMART de cada objetivo específico. El
criterio medible y el hito de cierre están resumidos en el
Cuadro~\ref{tab:smart} de la Sección~\ref{subsec:smart}.

\textbf{OE1.} \emph{Específico}: compara tres detectores concretos, STA/LTA,
EQTransformer y PhaseNet, sobre un conjunto de datos delimitado.
\emph{Medible}: sensibilidad (\emph{recall}) y tasa de falsos positivos,
globales y por clase, con intervalos de confianza del 95\,\% obtenidos por
\emph{bootstrap}; distribución de IoU; error de onset; frontera de Pareto de
los barridos de parámetros.
\emph{Alcanzable}: ejecutado en su totalidad con inferencia en CPU sobre
modelos preentrenados, sin requerir entrenamiento.
\emph{Relevante}: aborda los aspectos P1 y P2 del problema
(Cuadro~\ref{tab:trazabilidad}). \emph{Temporizado}: cerrado y documentado
para la Etapa 2, el 2 de octubre de 2026.

\textbf{OE2.} \emph{Específico}: define un procedimiento de generación de
espectrogramas cuya extensión de recorte depende de la duración detectada, y
que tolera los errores de delimitación temporal y el solapamiento entre
eventos. \emph{Medible}: IoU del recorte adaptativo contra el recorte de
referencia, comparado con la línea base de ventana fija y desglosado por
clase; fracción de eventos truncados; cobertura del rango de duraciones de
3,4\,s a 240,8\,s; número de solapamientos resueltos correctamente; tolerancia
al error de onset medida con desplazamientos controlados del inicio.
\emph{Alcanzable}: opera sobre las salidas ya
generadas por OE1, sin dependencia de hardware especializado.
\emph{Relevante}: aborda P3 y P1. \emph{Temporizado}: diseñado para la
Etapa 2 e implementado y medido para la Etapa 3, el 30 de octubre de 2026.

\textbf{OE3.} \emph{Específico}: ajusta por transferencia de aprendizaje una
única arquitectura convolucional, VGG16 preentrenada en ImageNet con capa de
salida de cuatro clases, sobre los espectrogramas producidos por OE2.
\emph{Medible}: F1 macro y \emph{recall} por clase con intervalos de confianza
del 95\,\%; matriz de confusión de cuatro clases; protocolo de entrenamiento y
evaluación reproducible, verificado por re-ejecución con la misma semilla.
\emph{Alcanzable}: el entrenamiento se realiza en Google Colab y la inferencia
en CPU local; el conjunto de entrenamiento de 9\,102 eventos etiquetados ya
está disponible, y ajustar una sola red reduce el presupuesto respecto de la
comparación de arquitecturas prevista originalmente. \emph{Relevante}: aborda
P5. \emph{Temporizado}: resultados preliminares en la Etapa 3 y ajuste
completo en la Etapa 4, el 4 de diciembre de 2026.

\textbf{OE4.} \emph{Específico}: evalúa la cadena completa sobre el registro
continuo, comparando el desempeño del clasificador en dos condiciones,
recorte de referencia y recorte automático, para los mismos eventos.
\emph{Medible}: diferencia de desempeño pareada entre ambas condiciones con su
intervalo de confianza; resultados desglosados por clase; tiempo de
procesamiento de la cadena completa en CPU. \emph{Alcanzable}: reutiliza los
artefactos de OE1 a OE3 y solo requiere inferencia, ejecutable en el equipo
local. \emph{Relevante}: aborda P4. \emph{Temporizado}: cerrado en la
Etapa 4, el 4 de diciembre de 2026.


\section{Evidencia empírica de la evaluación de detectores}
\label{sec:anexo-evidencia}

La evaluación de detectores, cerrada antes de esta etapa, entregó nueve
hallazgos que condicionan el diseño de los bloques posteriores. Todos
corresponden a mediciones sobre la traza continua de 10,18~horas con su catálogo
de 203 eventos, y son los que la Sección~\ref{sec:evidencia} invoca como origen
de los requerimientos funcionales.

{\small
\begin{longtable}{>{\raggedright\arraybackslash}p{0.9cm}>{\raggedright\arraybackslash}p{5.3cm}>{\raggedright\arraybackslash}p{5.3cm}}
\caption{Hallazgos medidos en la evaluación de detectores y su implicación de diseño.}\label{tab:evidencia}\\
\toprule
\textbf{ID} & \textbf{Hallazgo medido} & \textbf{Implicación de diseño} \\
\midrule
\endfirsthead
\multicolumn{3}{l}{\itshape Cuadro \thetable{} (continuación)}\\
\toprule
\textbf{ID} & \textbf{Hallazgo medido} & \textbf{Implicación de diseño} \\
\midrule
\endhead
\midrule
\multicolumn{3}{r}{\itshape Continúa en la página siguiente}\\
\endfoot
\bottomrule
\endlastfoot
E1 & El recall global de los tres detectores es estadísticamente indistinguible:
0,542; 0,527 y 0,581, con intervalos de confianza del 95\,\% solapados &
El detector no puede seleccionarse por su recall global \\
\addlinespace
E2 & Por clase existen diferencias significativas de signo opuesto que se cancelan
al promediar: un detector aventaja en LP y TR, el otro en VT, AV e IC &
Toda métrica debe reportarse desagregada por clase \\
\addlinespace
E3 & En calidad de delimitación, STA/LTA supera a EQTransformer de forma robusta
frente al umbral de solapamiento: IoU de 0,563 contra 0,479 &
La variabilidad del recorte entre detectores es un hecho medido; el clasificador
debe absorberla \\
\addlinespace
E4 & PhaseNet no delimita eventos: entrega arribos, no intervalos &
Se requiere una regla explícita de expansión de arribo a intervalo \\
\addlinespace
E5 & La mediana de duración de TR es de 87,8~s y excede la ventana fija de
81,92~s; el rango observado va de 3,4~s a 240,8~s &
El recorte debe ser adaptativo: la ventana fija trunca la clase más extensa \\
\addlinespace
E6 & 73 de los 203 eventos del catálogo continuo, un 36\,\%, son de clases no
volcánicas, y los detectores las encuentran con buen desempeño &
El rechazo no puede ocurrir en la detección; debe resolverse en la clasificación \\
\addlinespace
E7 & De 106 falsas alarmas analizadas, un 20,8\,\% son fragmentos de eventos
catalogados, un 35,8\,\% son adyacentes y un 43,4\,\% aisladas &
Las falsas alarmas son material de entrenamiento aprovechable como negativos
duros \\
\addlinespace
E8 & La estrategia de relleno de componentes altera el techo de recall y su efecto
depende del modelo: uno pasa de 0,527 a 0,946, el otro se comporta de forma
inversa &
El preprocesamiento es un parámetro del sistema, no una constante \\
\addlinespace
E9 & El error de onset se sitúa entre 1,7~s y 2,0~s en los modelos profundos y en
2,4~s en STA/LTA &
El recorte debe incorporar margen de contexto previo \\
\bottomrule
\end{longtable}
}


\section{Desarrollo de la evaluación de alternativas}
\label{sec:anexo-alternativas}

Se detallan aquí las siete decisiones de diseño resumidas en la
Sección~\ref{sec:alternativas}. El Cuadro~\ref{tab:alternativas} las presenta en
forma comparativa y los párrafos siguientes desarrollan el criterio decisivo en
cada caso.

{\small
\begin{longtable}{>{\raggedright\arraybackslash}p{0.8cm}>{\raggedright\arraybackslash}p{3.0cm}>{\raggedright\arraybackslash}p{3.6cm}>{\raggedright\arraybackslash}p{3.7cm}}
\caption{Alternativas consideradas y criterio de selección.}\label{tab:alternativas}\\
\toprule
\textbf{ID} & \textbf{Decisión} & \textbf{Alternativas} & \textbf{Selección y criterio} \\
\midrule
\endfirsthead
\multicolumn{4}{l}{\itshape Cuadro \thetable{} (continuación)}\\
\toprule
\textbf{ID} & \textbf{Decisión} & \textbf{Alternativas} & \textbf{Selección y criterio} \\
\midrule
\endhead
\midrule
\multicolumn{4}{r}{\itshape Continúa en la página siguiente}\\
\endfoot
\bottomrule
\endlastfoot
D1 & Elección del detector &
Detector único elegido por desempeño; conservar varios &
Conservar varios: el recorte variable es el objeto de estudio, no un defecto a
eliminar \\
\addlinespace
D2 & Manejo de clases no volcánicas &
Umbral en detección; detección OOD posterior; clase de rechazo supervisada &
Rechazo supervisado: hay 1.782 ejemplos etiquetados y los detectores sí
encuentran esas clases \\
\addlinespace
D3 & Estrategia de recorte &
Ventana fija de 81,92~s; ventana fija mayor; recorte adaptativo &
Recorte adaptativo: la ventana fija trunca TR y una mayor degrada la relación
entre evento y contexto \\
\addlinespace
D4 & Representación de entrada &
Forma de onda cruda; espectrograma; escalograma &
Representación espectral: mantiene comparabilidad con la línea base publicada \\
\addlinespace
D5 & Estrategia de entrenamiento del clasificador &
Entrenar desde cero; transferencia desde ImageNet; transferencia desde los
pesos de la línea base &
Transferencia desde ImageNet: mantiene comparable la arquitectura con la línea
base y cabe en el presupuesto \\
\addlinespace
D6 & Uso de detectores profundos &
Preentrenados; ajuste fino con datos locales &
Preentrenados: el ajuste fino consumiría el único conjunto con delimitación de
referencia \\
\addlinespace
D7 & Nivel de operación &
Canal individual; coincidencia multiestación &
Canal individual: la asociación robusta excede el presupuesto de esta etapa \\
\bottomrule
\end{longtable}
}

\textbf{D1. Elección del detector.} La alternativa natural era seleccionar el
mejor detector y descartar los demás. La medición lo impide: E1 muestra que sus
recalls globales son indistinguibles, y E2 muestra que las diferencias por clase
tienen signo opuesto, de modo que no existe un ganador único. Además, E3
establece que los recortes que producen difieren de forma medible. Conservar más
de un detector convierte esa variabilidad en la variable independiente del
estudio, que es exactamente lo que exige el objetivo general. Para la cadena
integrada se emplea como fuente principal el detector que mejor delimita, y se
mide también con el segundo, para cuantificar la sensibilidad del clasificador a
la fuente del recorte.

\textbf{D2. Manejo de las clases no volcánicas.} Se consideraron tres caminos.
Elevar el umbral de detección para que las avalanchas y los sismos de hielo no se
detecten sacrificaría recall, porque E6 muestra que los detectores las encuentran
bien y que constituyen el 36\,\% del catálogo continuo. Emplear un detector de
datos fuera de distribución posterior al clasificador, como hace el trabajo de
referencia \parencite{espinosacurilem2025ood}, exige calibración adicional y desaprovecha los 1.782 ejemplos etiquetados de
esas clases que están disponibles. Entrenar una clase de rechazo supervisada los
aprovecha y permite incorporar negativos duros extraídos de las falsas alarmas,
según E7. Se opta por la tercera y se conserva la segunda como línea base de
comparación. La limitación de esta elección debe declararse: un rechazo
supervisado no generaliza a clases nunca observadas, cosa que un detector de
datos fuera de distribución sí aspira a hacer.

\textbf{D3. Estrategia de recorte.} La ventana fija de 8.192 muestras trunca la
clase TR, cuya mediana de duración excede la ventana según E5. Ampliar la ventana
fija hasta cubrir el evento más largo del catálogo exigiría unas 24.080 muestras,
con lo que un evento de 3,4~s ocuparía cerca del 1,4\,\% de la ventana y el resto
sería contexto sin información de clase, además de multiplicar el costo de
cómputo. El recorte adaptativo preserva la proporción entre evento y contexto,
pero pierde la escala absoluta de duración al normalizar, por lo que esa
información debe reinyectarse de forma explícita en la representación.

\textbf{D4. Representación de entrada.} Las alternativas eran alimentar la red
con la forma de onda cruda o con una representación tiempo-frecuencia. Se opta
por la representación espectral porque el trabajo de referencia sobre este mismo
complejo volcánico la emplea \parencite{espinosacurilem2026segmentacion} y mantener esa elección permite atribuir las diferencias de desempeño al recorte
y no a la representación. Declaro que la comparación empírica entre
representaciones queda fuera del alcance por presupuesto de tiempo: es una
decisión de comparabilidad, no un resultado medido.

\textbf{D5. Estrategia de entrenamiento del clasificador.} La formulación
original de esta decisión comparaba tres arquitecturas convolucionales; la
reformulación de los objetivos en la retroalimentación de la Etapa 1 la
convierte en una decisión sobre cómo entrenar una sola. Se consideraron tres
caminos: entrenar la red desde cero, ajustarla por transferencia de
aprendizaje desde pesos preentrenados en ImageNet, o ajustarla desde los pesos
de la línea base publicada sobre este mismo complejo volcánico
\parencite{espinosacurilem2025ood}. Se adopta la transferencia desde ImageNet
sobre VGG16 con capa de salida de cuatro clases. La arquitectura coincide con
la de la línea base, lo que mantiene comparables los resultados y permite
atribuir las diferencias de desempeño al recorte y no al modelo; entrenar
desde cero exigiría un volumen de datos y un presupuesto de cómputo mayores
que los disponibles; y la disponibilidad de los pesos de la línea base no está
confirmada, por lo que esa alternativa se consultará al profesor guía y, de
obtenerse, se documentará como variante del mismo protocolo. Una sola VGG16
infiere en CPU sin dificultad, lo que atiende la restricción del entorno de
ejecución.

\textbf{D6. Uso de los detectores profundos.} Ajustar finamente PhaseNet \parencite{zhu2019phasenet} o EQTransformer \parencite{mousavi2020eqtransformer} con datos locales podría mejorar su desempeño, pero exigiría particionar el
catálogo continuo, que es el único conjunto con delimitación de referencia y está
reservado para evaluación. Se emplean preentrenados y se declara el ajuste fino
como trabajo futuro.

\textbf{D7. Nivel de operación.} La coincidencia entre estaciones reduce falsos
positivos, pero requiere una regla de asociación robusta cuya validación excede
el presupuesto de esta etapa. Se exploró de forma preliminar y se descartó para
el alcance actual. La detección opera sobre canales individuales y la asociación
multiestación queda como trabajo futuro.

\section{Detalle de los requerimientos}
\label{sec:anexo-rf}

Se listan los veinte requerimientos funcionales y los ocho no funcionales
sintetizados en la Sección~\ref{sec:sintesis-req}. La columna \emph{Origen}
remite a los hallazgos del Cuadro~\ref{tab:evidencia}, a los aspectos del
problema o a una restricción del entorno. La columna \emph{Verificación} indica
el procedimiento con el que se comprobará el cumplimiento.

{\small
\begin{longtable}{>{\raggedright\arraybackslash}p{1.2cm}>{\raggedright\arraybackslash}p{5.0cm}>{\raggedright\arraybackslash}p{1.6cm}>{\raggedright\arraybackslash}p{3.7cm}}
\caption{Requerimientos funcionales del bloque A, ingesta y preprocesamiento.}\label{tab:rf-a}\\
\toprule
\textbf{ID} & \textbf{Requerimiento} & \textbf{Origen} & \textbf{Verificación} \\
\midrule
\endfirsthead
\multicolumn{4}{l}{\itshape Cuadro \thetable{} (continuación)}\\
\toprule
\textbf{ID} & \textbf{Requerimiento} & \textbf{Origen} & \textbf{Verificación} \\
\midrule
\endhead
\midrule
\multicolumn{4}{r}{\itshape Continúa en la página siguiente}\\
\endfoot
\bottomrule
\endlastfoot
RF-01 & Leer la traza continua multicanal a 100~Hz, interpretando la primera fila
como marca de tiempo, las ocho siguientes como estaciones y las seis finales como
máscaras de clase & Datos &
Prueba de lectura que confirme dimensiones y rango temporal de los tres bloques \\
\addlinespace
RF-02 & Operar únicamente sobre los canales con señal utilizable y documentar la
exclusión de los restantes & Diagnóstico propio &
Informe de estadísticos por canal que sustente cada exclusión \\
\addlinespace
RF-03 & Construir la entrada de tres componentes que exigen los detectores
profundos mediante una estrategia de relleno declarada y configurable por modelo &
E8 & Reproducir los techos de recall medidos bajo cada estrategia \\
\addlinespace
RF-04 & Excluir las guardas de unión entre bloques temporales al construir el
conjunto de evaluación & Estructura del archivo &
Conteo por clase igual a 38 VT, 55 LP, 37 TR, 23 AV y 50 IC \\
\addlinespace
RF-05 & Garantizar una partición que impida que material de entrenamiento aparezca
en el conjunto de evaluación & RG-01 &
Prueba de correlación de forma de onda entre ambos conjuntos, documentada \\
\bottomrule
\end{longtable}
}

{\small
\begin{longtable}{>{\raggedright\arraybackslash}p{1.2cm}>{\raggedright\arraybackslash}p{5.0cm}>{\raggedright\arraybackslash}p{1.6cm}>{\raggedright\arraybackslash}p{3.7cm}}
\caption{Requerimientos funcionales del bloque B, detección.}\label{tab:rf-b}\\
\toprule
\textbf{ID} & \textbf{Requerimiento} & \textbf{Origen} & \textbf{Verificación} \\
\midrule
\endfirsthead
\multicolumn{4}{l}{\itshape Cuadro \thetable{} (continuación)}\\
\toprule
\textbf{ID} & \textbf{Requerimiento} & \textbf{Origen} & \textbf{Verificación} \\
\midrule
\endhead
\midrule
\multicolumn{4}{r}{\itshape Continúa en la página siguiente}\\
\endfoot
\bottomrule
\endlastfoot
RF-06 & Exponer una interfaz común de detección, de modo que cualquier detector
entregue una lista de intervalos con inicio, término y puntaje & E3, diseño &
Los tres detectores implementan la misma interfaz; prueba con un detector
sintético de salida conocida \\
\addlinespace
RF-07 & Aplicar una regla explícita de expansión de arribo a intervalo para los
detectores que no delimitan eventos & E4 &
Medición del IoU resultante tras aplicar la regla \\
\addlinespace
RF-08 & Reportar las métricas de detección desagregadas por clase, además de la
métrica global & E1, E2 &
Cada tabla de resultados incluye las cinco clases del catálogo \\
\bottomrule
\end{longtable}
}

{\small
\begin{longtable}{>{\raggedright\arraybackslash}p{1.2cm}>{\raggedright\arraybackslash}p{5.0cm}>{\raggedright\arraybackslash}p{1.6cm}>{\raggedright\arraybackslash}p{3.7cm}}
\caption{Requerimientos funcionales del bloque C, recorte y representación.}\label{tab:rf-c}\\
\toprule
\textbf{ID} & \textbf{Requerimiento} & \textbf{Origen} & \textbf{Verificación} \\
\midrule
\endfirsthead
\multicolumn{4}{l}{\itshape Cuadro \thetable{} (continuación)}\\
\toprule
\textbf{ID} & \textbf{Requerimiento} & \textbf{Origen} & \textbf{Verificación} \\
\midrule
\endhead
\midrule
\multicolumn{4}{r}{\itshape Continúa en la página siguiente}\\
\endfoot
\bottomrule
\endlastfoot
RF-09 & Recortar cada evento con una duración adaptada al intervalo entregado por
el detector, cubriendo el rango de duraciones observado & E5, P3 &
Distribución de duraciones de los recortes contrastada con la del catálogo \\
\addlinespace
RF-10 & Aplicar una política explícita para los eventos cuya duración excede la
ventana nominal, sin truncamiento arbitrario & E5, P3 &
Prueba sobre los eventos de mayor duración, con inspección visual \\
\addlinespace
RF-11 & Incorporar margen de contexto previo al inicio estimado para tolerar el
error de onset del detector & E9 &
Ensayo de sensibilidad frente a desplazamientos controlados del inicio \\
\addlinespace
RF-12 & Resolver el solapamiento entre detecciones consecutivas mediante una regla
documentada & OE2 &
Casos sintéticos con solapamiento conocido y resultado esperado \\
\addlinespace
RF-13 & Transformar recortes de duración variable en una representación espectral
de dimensiones fijas, preservando información de duración & E5, OE2 &
Comparación del desempeño con y sin preservación explícita de la duración \\
\bottomrule
\end{longtable}
}

{\small
\begin{longtable}{>{\raggedright\arraybackslash}p{1.2cm}>{\raggedright\arraybackslash}p{5.0cm}>{\raggedright\arraybackslash}p{1.6cm}>{\raggedright\arraybackslash}p{3.7cm}}
\caption{Requerimientos funcionales del bloque D, clasificación.}\label{tab:rf-d}\\
\toprule
\textbf{ID} & \textbf{Requerimiento} & \textbf{Origen} & \textbf{Verificación} \\
\midrule
\endfirsthead
\multicolumn{4}{l}{\itshape Cuadro \thetable{} (continuación)}\\
\toprule
\textbf{ID} & \textbf{Requerimiento} & \textbf{Origen} & \textbf{Verificación} \\
\midrule
\endhead
\midrule
\multicolumn{4}{r}{\itshape Continúa en la página siguiente}\\
\endfoot
\bottomrule
\endlastfoot
RF-14 & Entregar cuatro salidas: las tres clases volcánicas y una clase de
rechazo & E6, P5 &
Matriz de confusión de cuatro clases sobre el conjunto de evaluación \\
\addlinespace
RF-15 & Entrenar la clase de rechazo con eventos no volcánicos, ventanas de fondo
y negativos duros extraídos de las falsas alarmas & E7 &
Registro del origen y volumen de cada componente del conjunto de rechazo \\
\addlinespace
RF-16 & Aplicar y documentar un tratamiento explícito del desbalance entre clases &
Conteos del conjunto &
Reporte de métricas macro y por clase antes y después del tratamiento \\
\addlinespace
RF-17 & Ajustar por transferencia de aprendizaje una CNN preentrenada, bajo un
protocolo de entrenamiento y evaluación reproducible & OE3 &
La re-ejecución con la misma semilla reproduce las métricas dentro de la
tolerancia declarada \\
\bottomrule
\end{longtable}
}

{\small
\begin{longtable}{>{\raggedright\arraybackslash}p{1.2cm}>{\raggedright\arraybackslash}p{5.0cm}>{\raggedright\arraybackslash}p{1.6cm}>{\raggedright\arraybackslash}p{3.7cm}}
\caption{Requerimientos funcionales del bloque de evaluación integrada.}\label{tab:rf-e}\\
\toprule
\textbf{ID} & \textbf{Requerimiento} & \textbf{Origen} & \textbf{Verificación} \\
\midrule
\endfirsthead
\multicolumn{4}{l}{\itshape Cuadro \thetable{} (continuación)}\\
\toprule
\textbf{ID} & \textbf{Requerimiento} & \textbf{Origen} & \textbf{Verificación} \\
\midrule
\endhead
\midrule
\multicolumn{4}{r}{\itshape Continúa en la página siguiente}\\
\endfoot
\bottomrule
\endlastfoot
RF-18 & Permitir la evaluación pareada del clasificador sobre recorte de
referencia y sobre recorte automático, para los mismos eventos & P4, OE4 &
Ambas condiciones se ejecutan sobre el mismo conjunto de eventos \\
\addlinespace
RF-19 & Acompañar toda métrica comparativa de su intervalo de confianza estimado
por remuestreo & OE4 &
Cada comparación reporta su intervalo del 95\,\% \\
\addlinespace
RF-20 & Persistir resultados y configuración de cada ejecución de forma que
cualquier cifra reportada pueda regenerarse & OE4, RNF-03 &
Regeneración de al menos una tabla de resultados desde el registro almacenado \\
\bottomrule
\end{longtable}
}

{\small
\begin{longtable}{>{\raggedright\arraybackslash}p{1.4cm}>{\raggedright\arraybackslash}p{2.6cm}>{\raggedright\arraybackslash}p{7.4cm}}
\caption{Requerimientos no funcionales.}\label{tab:rnf}\\
\toprule
\textbf{ID} & \textbf{Característica} & \textbf{Requerimiento y verificación} \\
\midrule
\endfirsthead
\multicolumn{3}{l}{\itshape Cuadro \thetable{} (continuación)}\\
\toprule
\textbf{ID} & \textbf{Característica} & \textbf{Requerimiento y verificación} \\
\midrule
\endhead
\midrule
\multicolumn{3}{r}{\itshape Continúa en la página siguiente}\\
\endfoot
\bottomrule
\endlastfoot
RNF-01 & Portabilidad &
La inferencia debe ejecutarse íntegramente en CPU, sin dependencia de CUDA. Se
verifica ejecutando la cadena completa en el equipo local \\
\addlinespace
RNF-02 & Portabilidad &
El entrenamiento debe poder realizarse en un entorno de cuaderno remoto con sesión
limitada, con persistencia periódica de estado. Se verifica reanudando un
entrenamiento tras una interrupción \\
\addlinespace
RNF-03 & Fiabilidad &
Toda ejecución debe fijar semillas y registrar versiones de las bibliotecas. Se
verifica repitiendo una ejecución y comparando resultados \\
\addlinespace
RNF-04 & Mantenibilidad &
El detector debe ser intercambiable sin modificar el clasificador. Se verifica
sustituyéndolo y ejecutando la cadena sin cambios en el bloque D \\
\addlinespace
RNF-05 & Eficiencia &
El tiempo de procesamiento de la traza completa en CPU debe medirse y reportarse.
La meta de diseño es procesarla en menos tiempo que su duración real; el valor
alcanzado se reporta con independencia de si se cumple \\
\addlinespace
RNF-06 & Adecuación funcional &
Ninguna comparación entre configuraciones puede sustentarse en una diferencia
puntual sin medida de incertidumbre asociada \\
\addlinespace
RNF-07 & Mantenibilidad &
El código debe estar versionado en el repositorio, organizado por etapa,
excluyendo datos y pesos de modelos del control de versiones \\
\addlinespace
RNF-08 & Usabilidad &
Las salidas deben incluir un archivo tabular de detecciones clasificadas y figuras
de inspección visual, para permitir la revisión manual de casos \\
\bottomrule
\end{longtable}
}

\section{Criterios de aceptación}
\label{sec:anexo-aceptacion}

Los criterios se formulan como comparaciones contra una referencia medible y
no como umbrales absolutos, porque fijar un umbral numérico sin medición previa
que lo respalde sería arbitrario.

\begin{itemize}
\item \textbf{CA-01.} El clasificador entrenado sobre recortes adaptativos debe
superar, con intervalos de confianza que no se solapen, a un clasificador de
referencia entrenado sobre la ventana fija actual, en las métricas por clase.
\item \textbf{CA-02.} La degradación entre recorte de referencia y recorte
automático debe quedar cuantificada con su intervalo de confianza. No se fija un
máximo admisible: esa magnitud es el resultado que OE4 busca establecer.
\item \textbf{CA-03.} La clase de rechazo debe reducir la proporción de
detecciones no volcánicas propagadas como eventos volcánicos, respecto de la
configuración sin rechazo.
\item \textbf{CA-04.} El ajuste por transferencia debe reportarse con la
partición, la semilla y el presupuesto de entrenamiento declarados, y con las
métricas desglosadas por clase.
\item \textbf{CA-05.} La cadena completa debe ejecutarse sobre la traza de
evaluación en CPU y reportar su tiempo de procesamiento.
\end{itemize}

\section{Riesgos y mitigaciones}
\label{sec:anexo-riesgos}

Se listan los riesgos que pueden comprometer la validez de los resultados o el
cumplimiento del calendario. Los dos de severidad alta se enuncian además en la
Sección~\ref{sec:riesgos}. La severidad corresponde a un juicio propio y se
declara como tal, no como una estimación cuantitativa.

{\small
\begin{longtable}{>{\raggedright\arraybackslash}p{1.1cm}>{\raggedright\arraybackslash}p{4.4cm}>{\raggedright\arraybackslash}p{1.5cm}>{\raggedright\arraybackslash}p{4.5cm}}
\caption{Riesgos identificados y mitigación prevista.}\label{tab:riesgos}\\
\toprule
\textbf{ID} & \textbf{Riesgo} & \textbf{Severidad} & \textbf{Mitigación} \\
\midrule
\endfirsthead
\multicolumn{4}{l}{\itshape Cuadro \thetable{} (continuación)}\\
\toprule
\textbf{ID} & \textbf{Riesgo} & \textbf{Severidad} & \textbf{Mitigación} \\
\midrule
\endhead
\midrule
\multicolumn{4}{r}{\itshape Continúa en la página siguiente}\\
\endfoot
\bottomrule
\endlastfoot
RG-01 & Superposición entre el conjunto de entrenamiento y la traza de
evaluación. Hay 31 archivos cuya fecha parcial cae dentro de los bloques de la
traza, pero los nombres carecen de año y los conteos son compatibles con el azar.
No está probada ni descartada &
Alta &
Partición por bloque temporal, más verificación directa por correlación de forma
de onda antes de reportar resultados \\
\addlinespace
RG-02 & Discrepancia en el número de estaciones con señal entre el conjunto de
entrenamiento y la traza continua &
Media &
Auditoría sobre una muestra amplia y adecuación del preprocesamiento si se
confirma \\
\addlinespace
RG-03 & Truncamiento sistemático de los eventos más extensos por efecto de la
ventana fija &
Alta &
Recorte adaptativo con política explícita para eventos extensos \\
\addlinespace
RG-04 & Desbalance entre clases en el conjunto de entrenamiento, cercano a cuatro
a uno entre la más y la menos representada &
Media &
Tratamiento explícito del desbalance y reporte de métricas por clase \\
\addlinespace
RG-05 & Limitaciones de sesión del entorno remoto de entrenamiento &
Media &
Persistencia periódica de estado y experimentos fraccionables \\
\addlinespace
RG-06 & Cambios de alcance derivados de la retroalimentación del profesor guía &
Baja &
Objetivos formulados en términos de lo que se mide y no de cómo se implementa \\
\bottomrule
\end{longtable}
}

\section{Capacidad, método de estimación y épicas}
\label{sec:anexo-planificacion}

La capacidad de trabajo se estimó al inicio en 20~horas semanales, valor medio de
una disponibilidad declarada entre 15 y 25~horas, sobre la que se aplica un factor
de disponibilidad efectiva de 0,8 para absorber interrupciones y tareas
administrativas. Resultaban 16~horas comprometibles por semana y una capacidad
total cercana a 208~horas para las 13,1 semanas que restan hasta la entrega final,
repartidas en tres sprints de cuatro semanas.

Esa cifra dejó de ser válida antes de cerrar esta etapa. Inicié mi práctica
profesional, mayormente remota y con asistencia presencial los lunes, lo que sitúa
la disponibilidad real cerca de 15~horas semanales. Con el mismo factor de 0,8
resultan 12~horas comprometibles por semana y una capacidad revisada cercana a
157~horas, un 25\,\% menos que la estimación inicial.

El backlog suma 181~horas y por lo tanto excede esa capacidad revisada. En lugar
de recortar tareas de forma prematura, traslado el ajuste a la prioridad: las 30
tareas se reparten en 142~horas de prioridad \emph{Debe}, 36 de \emph{Debería} y 3
de \emph{Podría}. Las 142~horas de prioridad \emph{Debe} caben dentro de las
157~disponibles y constituyen el compromiso firme de esta planificación; las de
prioridad \emph{Debería} y \emph{Podría} quedan declaradas como alcance opcional,
a ejecutar solo si la disponibilidad real supera lo previsto. Las cifras de
capacidad que encabezan los cuadros de cada sprint corresponden a la estimación
inicial y se conservan como referencia contra la cual medir la desviación.

Las estimaciones por tarea provienen de juicio propio anclado en trabajo análogo
ya ejecutado durante la evaluación de detectores. Corresponde declarar la
limitación: no se llevó registro del esfuerzo real invertido en aquellas tareas,
por lo que el anclaje es cualitativo y las cifras son estimaciones, no
mediciones. Como mitigación, registraré el esfuerzo real de cada tarea a medida
que se complete, de modo que los sprints siguientes puedan recalibrarse con datos
propios.

El backlog se organiza en seis épicas. \textbf{EP-1, validez experimental},
garantiza que ningún resultado esté contaminado por superposición entre conjuntos
ni por supuestos no verificados. \textbf{EP-2, cadena de detección operativa},
permite conectar cualquier detector mediante una interfaz única. \textbf{EP-3,
recorte adaptativo y representación}, resuelve el paso de duración variable a
entrada fija y responde a OE2. \textbf{EP-4, clasificador de cuatro salidas},
distingue las clases volcánicas y rechaza el resto, y responde a OE3.
\textbf{EP-5, evaluación
integrada}, mide la degradación entre recorte de referencia y automático, y
responde a OE4. \textbf{EP-6, documentación y cierre}, asegura que cada resultado
quede escrito, versionado y reproducible.

Cada épica se descompone en historias de usuario y estas en tareas técnicas
estimadas y priorizadas, detalladas en la Sección~\ref{sec:anexo-backlog}.


\section{Backlog del proyecto}
\label{sec:anexo-backlog}

El backlog se estructura en tres niveles: épicas, presentadas en la
Sección~\ref{sec:anexo-planificacion}; historias de usuario, que expresan la necesidad desde el
punto de vista de quien recibe el valor; y tareas técnicas, que son las unidades
estimadas y ejecutables.

\subsection{Historias de usuario}
\label{sec:anexo-historias}

{\small
\begin{longtable}{>{\raggedright\arraybackslash}p{1.2cm}>{\raggedright\arraybackslash}p{7.6cm}>{\raggedright\arraybackslash}p{1.0cm}>{\raggedright\arraybackslash}p{1.9cm}}
\caption{Historias de usuario, con su épica y las tareas que las realizan.}\label{tab:historias}\\
\toprule
\textbf{ID} & \textbf{Historia} & \textbf{Épica} & \textbf{Tareas} \\
\midrule
\endfirsthead
\multicolumn{4}{l}{\itshape Cuadro \thetable{} (continuación)}\\
\toprule
\textbf{ID} & \textbf{Historia} & \textbf{Épica} & \textbf{Tareas} \\
\midrule
\endhead
\midrule
\multicolumn{4}{r}{\itshape Continúa en la página siguiente}\\
\endfoot
\bottomrule
\endlastfoot
HU-01 & Como responsable del estudio, quiero garantizar que ningún evento usado
para entrenar aparezca en la traza de evaluación, para que las métricas midan
generalización y no memorización & EP-1 & B-01, B-03 \\
\addlinespace
HU-02 & Como responsable del estudio, quiero conocer qué estaciones tienen señal
utilizable en cada conjunto, para que el preprocesamiento sea equivalente entre
entrenamiento y evaluación & EP-1 & B-02 \\
\addlinespace
HU-03 & Como responsable del estudio, quiero poder sustituir el detector sin
modificar el resto de la cadena, para medir su efecto sobre la clasificación
final & EP-2 & B-06, B-07 \\
\addlinespace
HU-04 & Como analista del observatorio, quiero que un detector que solo entrega
arribos produzca igualmente intervalos, para comparar todos los detectores bajo un
mismo criterio & EP-2 & B-11 \\
\addlinespace
HU-05 & Como analista del observatorio, quiero que los eventos largos se recorten
completos, para que el tremor no se evalúe truncado & EP-3 & B-08, B-12 \\
\addlinespace
HU-06 & Como responsable del estudio, quiero que la representación de entrada
conserve la duración del evento, para no perder un rasgo discriminante entre
clases & EP-3 & B-14, B-15 \\
\addlinespace
HU-07 & Como responsable del estudio, quiero saber cuánto afecta el error de onset
a la clasificación, para dimensionar el margen de contexto del recorte & EP-3 &
B-13, B-21 \\
\addlinespace
HU-08 & Como analista del observatorio, quiero que las avalanchas y los sismos de
hielo no se reporten como eventos volcánicos, para no revisar manualmente falsas
identificaciones & EP-4 & B-16, B-17, B-27 \\
\addlinespace
HU-09 & Como responsable del estudio, quiero ajustar una CNN preentrenada bajo
un protocolo reproducible, para que el resultado sea replicable & EP-4 & B-19,
B-20, B-22, B-23 \\
\addlinespace
HU-10 & Como analista del observatorio, quiero saber cuánto se degrada el
desempeño al usar recortes automáticos en vez de manuales, para decidir si la
cadena es utilizable & EP-5 & B-25, B-26 \\
\addlinespace
HU-11 & Como analista del observatorio, quiero un archivo de detecciones
clasificadas y figuras de inspección, para revisar los casos dudosos & EP-5 &
B-28 \\
\addlinespace
HU-12 & Como evaluador del trabajo, quiero poder regenerar cualquier cifra del
informe desde el repositorio, para verificar los resultados & EP-6 & B-09,
B-29 \\
\bottomrule
\end{longtable}
}

\subsection{Definición de terminado}
\label{sec:anexo-dod}

Una tarea se considera terminada cuando cumple las cuatro condiciones siguientes.
El código está en el repositorio, con mensaje de commit descriptivo y sin datos ni
pesos incluidos. Cuando la tarea implica una regla de decisión, esta fue probada
con al menos un caso sintético de solución conocida antes de aplicarse a datos
reales. Las cifras producidas quedan almacenadas junto con la configuración que
las generó. Y todo supuesto o limitación detectado durante la ejecución queda
registrado en el documento, no solo en el código.

\subsection{Sprint 1, cierre el 2 de octubre}
\label{sec:anexo-s1}

Concentra las condiciones habilitantes y el diseño de la cadena. La tarea B-01
es bloqueante para el entrenamiento: mientras no se resuelva la posible
superposición entre conjuntos, cualquier entrenamiento arriesga producir
resultados no interpretables. Como el entrenamiento es materia de la Etapa 3 y
las dos primeras semanas del sprint quedaron absorbidas por el inicio de la
práctica profesional, B-01 se posterga a la primera semana del sprint 2; la
desviación se declara aquí en lugar de absorberse en silencio. El foco de la
Etapa 2 son las cuatro tareas de tipo diseño del backlog, B-04, B-05, B-06 y
B-08. La columna
\emph{Tipo} distingue diseño, implementación, validación y documentación.

{\small
\begin{longtable}{>{\raggedright\arraybackslash}p{0.9cm}>{\raggedright\arraybackslash}p{4.5cm}>{\raggedright\arraybackslash}p{1.0cm}>{\raggedright\arraybackslash}p{1.9cm}>{\raggedright\arraybackslash}p{0.5cm}>{\raggedright\arraybackslash}p{1.2cm}}
\caption{Backlog del sprint 1. Capacidad comprometible según la estimación inicial: 64~horas.}\label{tab:bl1}\\
\toprule
\textbf{ID} & \textbf{Tarea} & \textbf{Tipo} & \textbf{Traza} & \textbf{h} & \textbf{Prior.} \\
\midrule
\endfirsthead
\multicolumn{6}{l}{\itshape Cuadro \thetable{} (continuación)}\\
\toprule
\textbf{ID} & \textbf{Tarea} & \textbf{Tipo} & \textbf{Traza} & \textbf{h} & \textbf{Prior.} \\
\midrule
\endhead
\midrule
\multicolumn{6}{r}{\itshape Continúa en la página siguiente}\\
\endfoot
\bottomrule
\endlastfoot
B-01 & Verificar por correlación de forma de onda si existe superposición entre el
conjunto de entrenamiento y la traza de evaluación & Val. & RF-05, RG-01 & 10 & Debe \\
\addlinespace
B-02 & Auditar sobre una muestra amplia el número de estaciones con señal en el
conjunto de entrenamiento & Val. & RG-02 & 5 & Debería \\
\addlinespace
B-03 & Implementar la partición por bloque temporal y documentar su composición
por clase & Imp. & RF-04, RF-05 & 4 & Debe \\
\addlinespace
B-04 & Elaborar el diagrama de arquitectura de la cadena de procesamiento, con sus
cuatro bloques y sus responsabilidades & Dis. & OE1, OE2 & 6 & Debe \\
\addlinespace
B-05 & Elaborar el diagrama de flujo del procesamiento y el contrato de datos
entre bloques & Dis. & RF-06 & 6 & Debe \\
\addlinespace
B-06 & Especificar la interfaz común de detección y su formato de intercambio &
Dis. & RF-06 & 5 & Debe \\
\addlinespace
B-07 & Implementar la interfaz común y validarla con un detector sintético de
salida conocida & Imp. & RF-06 & 7 & Debe \\
\addlinespace
B-08 & Diseñar el procedimiento de recorte adaptativo, con la política para
eventos que exceden la ventana nominal & Dis. & RF-09, RF-10 & 6 & Debe \\
\addlinespace
B-09 & Ordenar el repositorio por etapa, definir \texttt{.gitignore} y la
convención de mensajes de commit & Doc. & RNF-07 & 4 & Debe \\
\addlinespace
B-10 & Redactar el avance documental del sprint & Doc. & --- & 5 & Debe \\
\midrule
\multicolumn{4}{r}{\textbf{Comprometido}} & \textbf{58} & \\
\multicolumn{4}{r}{\textbf{Holgura}} & 6 & \\
\bottomrule
\end{longtable}
}

\subsection{Sprint 2, cierre el 30 de octubre}
\label{sec:anexo-s2}

Materializa el diseño del sprint anterior y produce los primeros resultados
preliminares del clasificador.

{\small
\begin{longtable}{>{\raggedright\arraybackslash}p{0.9cm}>{\raggedright\arraybackslash}p{4.5cm}>{\raggedright\arraybackslash}p{1.0cm}>{\raggedright\arraybackslash}p{1.9cm}>{\raggedright\arraybackslash}p{0.5cm}>{\raggedright\arraybackslash}p{1.2cm}}
\caption{Backlog del sprint 2. Capacidad comprometible según la estimación inicial: 64~horas.}\label{tab:bl2}\\
\toprule
\textbf{ID} & \textbf{Tarea} & \textbf{Tipo} & \textbf{Traza} & \textbf{h} & \textbf{Prior.} \\
\midrule
\endfirsthead
\multicolumn{6}{l}{\itshape Cuadro \thetable{} (continuación)}\\
\toprule
\textbf{ID} & \textbf{Tarea} & \textbf{Tipo} & \textbf{Traza} & \textbf{h} & \textbf{Prior.} \\
\midrule
\endhead
\midrule
\multicolumn{6}{r}{\itshape Continúa en la página siguiente}\\
\endfoot
\bottomrule
\endlastfoot
B-11 & Implementar la regla de expansión de arribo a intervalo y medir el IoU
resultante & Imp. & RF-07 & 7 & Debería \\
\addlinespace
B-12 & Implementar el recorte adaptativo según el diseño de B-08 & Imp. & RF-09,
RF-10 & 10 & Debe \\
\addlinespace
B-13 & Validar la regla de solapamiento con casos sintéticos de resultado
esperado & Val. & RF-12 & 4 & Debe \\
\addlinespace
B-14 & Implementar la representación espectral de dimensiones fijas a partir de
recortes de duración variable & Imp. & RF-13 & 10 & Debe \\
\addlinespace
B-15 & Incorporar la duración como información explícita y medir su aporte frente
a la representación sin ella & Imp. & RF-13 & 5 & Debería \\
\addlinespace
B-16 & Extraer el conjunto de negativos duros desde las falsas alarmas
caracterizadas & Imp. & RF-15, E7 & 4 & Debería \\
\addlinespace
B-17 & Construir el conjunto de entrenamiento de cuatro clases, incluida la clase
de rechazo & Imp. & RF-14, RF-15 & 7 & Debe \\
\addlinespace
B-18 & Redactar el avance documental del sprint & Doc. & --- & 5 & Debe \\
\addlinespace
B-19 & Entrenar el clasificador de línea base sobre ventana fija & Imp. & CA-01 &
6 & Debe \\
\midrule
\multicolumn{4}{r}{\textbf{Comprometido}} & \textbf{58} & \\
\multicolumn{4}{r}{\textbf{Holgura}} & 6 & \\
\bottomrule
\end{longtable}
}

\subsection{Sprint 3, cierre el 4 de diciembre}
\label{sec:anexo-s3}

Cierra OE3 y OE4 y produce el informe final. Es el más extenso y concentra el
trabajo de escritura, por lo que su holgura es proporcionalmente mayor. Tras la
reformulación de los objetivos, las tareas B-23 y B-24 de la Etapa 1 se funden
en una sola tarea B-23 de ajuste por transferencia; el identificador B-24 se
retira del backlog y no se reutiliza, para conservar la trazabilidad con la
numeración de la Etapa 1.

{\small
\begin{longtable}{>{\raggedright\arraybackslash}p{0.9cm}>{\raggedright\arraybackslash}p{4.5cm}>{\raggedright\arraybackslash}p{1.0cm}>{\raggedright\arraybackslash}p{1.9cm}>{\raggedright\arraybackslash}p{0.5cm}>{\raggedright\arraybackslash}p{1.2cm}}
\caption{Backlog del sprint 3. Capacidad comprometible según la estimación inicial: 80~horas.}\label{tab:bl3}\\
\toprule
\textbf{ID} & \textbf{Tarea} & \textbf{Tipo} & \textbf{Traza} & \textbf{h} & \textbf{Prior.} \\
\midrule
\endfirsthead
\multicolumn{6}{l}{\itshape Cuadro \thetable{} (continuación)}\\
\toprule
\textbf{ID} & \textbf{Tarea} & \textbf{Tipo} & \textbf{Traza} & \textbf{h} & \textbf{Prior.} \\
\midrule
\endhead
\midrule
\multicolumn{6}{r}{\itshape Continúa en la página siguiente}\\
\endfoot
\bottomrule
\endlastfoot
B-20 & Entrenar el clasificador sobre recorte adaptativo y compararlo con la línea
base & Imp. & CA-01 & 8 & Debe \\
\addlinespace
B-21 & Medir la sensibilidad del clasificador a desplazamientos controlados del
instante de inicio & Val. & RF-11, E9 & 4 & Debería \\
\addlinespace
B-22 & Aplicar el tratamiento del desbalance y medir su efecto por clase & Imp. &
RF-16, RG-04 & 4 & Debería \\
\addlinespace
B-23 & Ajustar VGG16 preentrenada en ImageNet a cuatro clases bajo el protocolo
común de entrenamiento y evaluación & Imp. & RF-17, CA-04 & 10 & Debe \\
\addlinespace
B-25 & Ejecutar la cadena completa sobre la traza de evaluación en CPU y medir su
tiempo de procesamiento & Val. & RNF-05, CA-05, OE4 & 6 & Debe \\
\addlinespace
B-26 & Realizar la evaluación pareada entre recorte de referencia y automático, con
intervalos de confianza & Val. & RF-18, RF-19 & 9 & Debe \\
\addlinespace
B-27 & Medir la reducción de detecciones no volcánicas propagadas que aporta la
clase de rechazo & Val. & CA-03 & 4 & Debería \\
\addlinespace
B-28 & Generar el archivo tabular de salida y las figuras de inspección visual &
Imp. & RNF-08 & 3 & Debería \\
\addlinespace
B-29 & Verificar la regeneración de una tabla de resultados desde el registro
almacenado & Val. & RF-20, RNF-03 & 3 & Podría \\
\addlinespace
B-30 & Redactar el capítulo de resultados & Doc. & --- & 6 & Debe \\
\addlinespace
B-31 & Redactar validación, conclusiones y recomendaciones para trabajo futuro &
Doc. & --- & 8 & Debe \\
\midrule
\multicolumn{4}{r}{\textbf{Comprometido}} & \textbf{65} & \\
\multicolumn{4}{r}{\textbf{Holgura}} & 15 & \\
\bottomrule
\end{longtable}
}

\subsection{Refinamiento de las tareas prioritarias}
\label{sec:anexo-refinamiento}

Las tareas de prioridad \emph{Debe} del primer sprint se refinan con criterios de
aceptación propios, adicionales a la definición de terminado.

\begin{itemize}
\item \textbf{B-01.} Se acepta cuando exista un informe que reporte, para cada uno
de los 31 archivos sospechosos, el coeficiente de correlación máximo contra la
traza y la decisión asociada, junto con la distribución de correlación de una
muestra de control de archivos no sospechosos.
\item \textbf{B-03.} Se acepta cuando la partición esté implementada como función
reproducible, y exista una tabla con el número de eventos por clase y por bloque
en cada subconjunto.
\item \textbf{B-04.} Se acepta cuando el diagrama identifique los cuatro bloques,
sus interfaces y el punto exacto donde se aplica cada requerimiento funcional, y
siga una notación estándar reconocida.
\item \textbf{B-05.} Se acepta cuando el contrato de datos entre bloques
especifique nombres de campo, tipos y unidades, y sea suficiente para que un
tercero implemente un bloque sin leer el código del anterior.
\item \textbf{B-06.} Se acepta cuando la especificación describa la interfaz de
forma independiente de cualquier detector concreto, y se demuestre que los tres
detectores evaluados encajan en ella.
\item \textbf{B-07.} Se acepta cuando el detector sintético, cuya salida es
conocida de antemano, atraviese la cadena y produzca exactamente los intervalos
esperados.
\item \textbf{B-08.} Se acepta cuando el diseño especifique el comportamiento para
los tres casos límite: evento más corto del catálogo, evento más largo y evento de
duración igual a la ventana nominal.
\item \textbf{B-09.} Se acepta cuando el repositorio no contenga datos ni pesos, y
el historial de commits permita identificar qué cambió en cada uno.
\end{itemize}

\subsection{Priorización, dependencias y política ante desviación}
\label{sec:anexo-priorizacion}

La prioridad se asigna según tres criterios ordenados: si la omisión de la tarea
invalidaría los resultados; si bloquea a otras tareas; y su aporte a la respuesta
de un objetivo específico. Las tareas marcadas como \emph{Podría} son las primeras
candidatas a salir del alcance si la capacidad real resulta inferior a la
estimada. Si la disponibilidad cae al extremo inferior del rango declarado, la
capacidad total baja cerca de un 25\,\% y la holgura de 27~horas no alcanza a
absorberla, por lo que la reducción debe hacerse por alcance y no por calidad.

Las dependencias que condicionan el orden de ejecución son las siguientes. B-01
bloquea a B-19, B-20 y B-23, porque entrenar antes de resolver la posible
superposición produciría métricas no interpretables. B-06 bloquea a B-07. B-08
bloquea a B-12, y B-12 bloquea a B-14, porque la representación se construye
sobre el recorte. B-17 bloquea a B-19, B-20 y B-23. B-25 bloquea a B-26.

Ante una desviación, si una tarea supera su estimación en más del 50\,\% se
detiene su ejecución y se reevalúa el alcance, en lugar de absorber la holgura de
forma silenciosa. Quedan explícitamente fuera del backlog las tareas
correspondientes a lo declarado fuera de alcance en la
Sección~\ref{sec:alcances}. No se incorporan como tareas de baja prioridad,
porque hacerlo sugeriría que podrían abordarse si sobrara tiempo, y no es el caso.